\section{胶子自能图}

\begin{figure}[b]
\centerline{\includegraphics[width=1.5in]{images/gluonselfl} \hspace{-5mm} \includegraphics[width=1.5in]{images/gluontriple}}
\caption{插入右腿的胶子自能型修正。
\label{gluself}}
\end{figure}

人们可以预期，
加入胶子自能图（其一示于图 \ref{gluself}a）之后
会出现加 prescription 形式。
这些图兼有 UV 与共线发散。在格点上，UV 发散由格距正规化。
在连续理论中，可对胶子传播子采用
Polyakov 规定 $1/z^2\to 1/(z^2 -a^2)$。
共线发散可通过取有限胶子质量 $\lambda$ 来正规化。结果是
$\ln (a^2 \lambda^2)$ 贡献。
但它不含 $z$ 依赖，
显然无助于为盒图贡献的
$\ln z_3^2$ 部分构建加 prescription 形式。

一种可能的出路是把 $\ln (a^2 \lambda^2)$
表示为演化型对数 $\ln (z_3^2 \lambda^2)$
与 UV 型对数 $\ln (z_3^2/a^2 )$ 之差 $\ln (z_3^2 \lambda^2)   -  \ln (z_3^2/a^2  )$。
后者可以并入分别对应于链接自能与顶点修正的
图 \ref{linkself} 与 \ref{link}
的 UV 发散之中。
而 $\ln (z_3^2 \lambda^2)$
部分则加入演化核。

为确保规范不变性起见，我们使用
维数正规化。
此时 $\ln (a^2 \lambda^2)$ 对数的类比是极点 $1/(2-d/2)$，有时写作 \mbox{$1/\epsilon_{\rm UV} - 1/\epsilon_{\rm IR}$。}
就我们的目的而言，更方便的是把它符号化地写成类似 $\ln (a^2 \lambda^2)$ 的形式。
作替换 $\lambda \to\mu_{\rm IR}$、$a  \to 1/\mu_{\rm UV}$ 得 $\ln (\mu_{\rm IR}^2/\mu_{\rm UV}^2)$，
然后把它拆分为差 $\ln (z_3^2 \mu_{\rm IR}^2)   -  \ln (z_3^2 \mu_{\rm UV}^2 )$。

还应计入从链接--胶子顶点额外伸出一条
胶子线的图（其一示于图 \ref{gluself}b）。图 \ref{gluself} 各图与其左腿类似图的合并贡献为
\begin{align}
{g^2N_c\over 8\pi^2}
\frac1{2-d/2}
\left [2 - \frac{\beta_0}{2N_c} \right ]G_{\mu\alpha}(z)G_{\lambda \beta }(0) \ ,
\label{counter3}
\end{align}
其中纯胶动力学中 $\beta_0 =11N_c/3$，故方括号内各项合成为 1/6。
如上所述，我们将把 \mbox{$1/({2-d/2})$} 视为
\mbox{$\ln (z_3^2 \mu_{\rm IR}^2)   -  \ln (z_3^2 \mu_{\rm UV}^2 )$。}
